% =============================================================
% Capitulo 4 - Diseno del circuito de deteccion de picos
% =============================================================
\chapter{Diseño del circuito de detección de picos}
\label{cap:diseno_deteccion_picos}

\section{Modelado del circuito de detección}
\label{sec:modelado_simulacion}

Para determinar el valor máximo de la señal de salida del fotodiodo se plantearon dos alternativas:

\begin{itemize}
  \item digitalizar la señal completa mediante un ADC rápido y
        procesarla posteriormente para localizar el valor máximo, o
  \item retener el valor máximo de la señal analógica con un circuito
        detector de picos y digitalizar ese valor con un ADC de menor
        velocidad.
\end{itemize}

Por razones de costo se adoptó la segunda alternativa, dado que admite un ADC más lento y económico y reduce el volumen de datos a procesar.

El circuito detector de picos diseñado consta de dos amplificadores realimentados por corriente (CFA), conectados como amplificadores no inversores, y un transistor bipolar que opera como elemento de detección. El amplificador X1 adapta la impedancia de salida del amplificador de transimpedancia de los fotodiodos. El transistor T2 se emplea como diodo detector en su unión base-emisor y carga el circuito RC conformado por $R_4$ y $C_1$. El amplificador X2 actúa como buffer de salida, evitando que la impedancia de entrada del ADC, representada mediante $R_7$, descargue el capacitor de retención antes de la digitalización. El empleo de CFA en lugar de los amplificadores realimentados por tensión (VFA) tradicionales permite alcanzar un mayor \emph{slew rate}, característica determinante paraseguir la rápida variación de la señal de salida del fotodiodo.

PLACEHOLDER PONER IMAGEN DE CIRCUITO DETECTOR DE PICOS

El amplificador X1 se configuró con una ganancia de 2 para aumentar el rango dinámico del circuito detector, puesto que con niveles bajos de tension caemos en la alinealidad del transisotr, dado por el diodo base-emisor.

Partiendo de la ecuación estándar de carga de un circuito RC,

\begin{equation}
  V_C(t) = V_O \left(1 - e^{-t/R_4 C_1}\right)
  \label{eq:carga_rc}
\end{equation}

\noindent y considerando que, para $t/R_4C_1 \ll 1$, dicha expresión
puede aproximarse como

\begin{equation}
  V_C(t) \approx V_O \, \frac{t}{R_4 C_1},
  \label{eq:carga_rc_lineal}
\end{equation}

\noindent se obtiene una relación lineal entre la tensión retenida y la tensión de pico de entrada sin necesidad de cargar por completo el capacitor $C_1$. Esto permite seleccionar valores de capacidad que otorguen tiempos de descarga suficientemente grandes para digitalizar el valor de pico $V_O$ con un conversor analógico-digital de velocidad moderada y bajo costo. El amplificador X2 se configuró con ganancia unitaria para evitar que la impedancia de entrada del ADC descargue el capacitor de retención antes de la digitalización.


\title{Selección de componentes del detector de picos}

El amplificador realimentado por corriente AD8004 se seleccionó por su elevado \emph{slew rate} y bajo costo, adecuado para la aplicación. El circuito integrado empleado en el prototipo fue provisto como muestra gratuita por Analog Devices. 

\begin{table}[h]
  \centering
  \begin{tabular}{lll}
    \toprule
    \textbf{Parámetro} & \textbf{Valor} & \textbf{Condiciones} \\
    \midrule
    \multicolumn{3}{l}{\textit{Desempeño dinámico}} \\
    Ancho de banda ($-3$~dB)            & 250~MHz            & $G = +1$          \\
    Ancho de banda (planitud 0.1~dB)    & 30~MHz             & $G = +2$          \\
    Slew rate                           & 3000~V/\textmu s   & $V_S = \pm5$~V    \\
    Slew rate                           & 1100~V/\textmu s   & $V_S = +5$~V      \\
    Tiempo de establecimiento (0.1\,\%) & 21~ns              &                   \\
    Tiempo de subida (escalón de 2~V)   & 1.8~ns             &                   \\
    \midrule
    \multicolumn{3}{l}{\textit{Ruido y distorsión}} \\
    Distorsión armónica total (THD)     & $-78$~dBc                    & $f = 5$~MHz  \\
    Ruido de tensión de entrada         & 1.5~nV/$\sqrt{\text{Hz}}$    & $f = 10$~kHz \\
    Ruido de corriente de entrada       & 38~pA/$\sqrt{\text{Hz}}$     & $f = 10$~kHz \\
    \midrule
    \multicolumn{3}{l}{\textit{Desempeño DC y características de entrada}} \\
    Tensión de \emph{offset} de entrada     & 1.0~mV (típ.); 3.5~mV (máx.) &                \\
    Corriente de polarización de entrada    & $\pm40$~\textmu A (típ.)      &                \\
    Transresistencia en lazo abierto        & 290~k$\Omega$                & $V_O = \pm2.5$~V \\
    Resistencia de entrada (no inversora)   & 2~M$\Omega$                  &                \\
    Rango de tensión de entrada (modo común)& 3.2~V                        &                \\
    \midrule
    \multicolumn{3}{l}{\textit{Salida y alimentación}} \\
    Corriente de salida                     & 50~mA               &                        \\
    Corriente de alimentación por amplificador & 3.5~mA (35~mW)   &                        \\
    Rango de alimentación (simple)          & +4 a +12~V          &                        \\
    Rango de alimentación (dual)            & $\pm2$ a $\pm6$~V   &                        \\
    Tensión de alimentación máxima          & 12.6~V              & Valor máximo absoluto  \\
    Rango de temperatura de operación       & $-40$ a $+85$~$^{\circ}$C & Grado A          \\
    Encapsulado                             & SOIC de 14 pines    &                        \\
    \bottomrule
  \end{tabular}
  \caption{Características principales del amplificador realimentado por
  corriente AD8004.}
  \label{tab:caracteristicas_ad8004}
\end{table}

Como elemento detector se utilizó un transistor bipolar BC817 debido a contar con alta corriente de colector, necesaria para cargar rápidamente el capacitor de emisor y además de poseer un ancho de banda lo suficientemente alto como para detectar los picos de la señal.

\begin{table}[htbp]
  \centering
  \begin{tabular}{lll}
    \toprule
    \textbf{Parámetro} & \textbf{Valor} & \textbf{Condiciones} \\
    \midrule
    \multicolumn{3}{l}{\textit{Valores máximos absolutos}} \\
    Tensión colector-base ($V_{CBO}$)   & 50~V             & $I_E = 0$        \\
    Tensión colector-emisor ($V_{CEO}$) & 45~V             & $I_B = 0$        \\
    Tensión emisor-base ($V_{EBO}$)     & 5~V              & $I_C = 0$        \\
    Corriente de colector ($I_C$)       & 500~mA           & Continua         \\
    Potencia de disipación ($P_D$)      & 300~mW           & $T_A = 25$~$^{\circ}$C \\
    Temperatura de unión ($T_J$)        & 150~$^{\circ}$C  & Máxima           \\
    \midrule
    \multicolumn{3}{l}{\textit{Características eléctricas}} \\
    Ganancia de corriente DC ($h_{FE}$)      & 160 a 400            & $V_{CE}=1$~V, $I_C=100$~mA \\
    Tensión de saturación C-E ($V_{CE(sat)}$)& 0.7~V (máx.)         & $I_C=500$~mA, $I_B=50$~mA  \\
    Tensión de saturación B-E ($V_{BE(sat)}$)& 1.2~V (máx.)         & $I_C=500$~mA, $I_B=50$~mA  \\
    Frecuencia de transición ($f_T$)         & 100~MHz (mín.)       & $V_{CE}=5$~V, $I_C=10$~mA  \\
    Corriente de fuga de colector ($I_{CBO}$)& 0.1~\textmu A (máx.) & $V_{CB}=45$~V             \\
    Corriente de fuga de emisor ($I_{EBO}$)  & 0.1~\textmu A (máx.) & $V_{EB}=4$~V              \\
    \midrule
    \multicolumn{3}{l}{\textit{Datos generales}} \\
    Tipo        & NPN de silicio, propósito general & \\
    Encapsulado & SOT-23                            & \\
    \bottomrule
  \end{tabular}
  \caption{Características principales del transistor bipolar NPN BC817.}
  \label{tab:caracteristicas_bc817}
\end{table}

Para la selección de los valores de los componentes pasivos se consideraron los siguientes criterios:
\begin{itemize}
      \item El valor de $R_2$ debe ser lo suficientemente grande para que el tiempo de carga del capacitor $C_1$ se mayor .
      \item Los valores de resistencias de realimentacion para los amplificadores U1A y U1B se seleccionaron para obtener una ganancia de 2 y 1 utilizando los valores recomendados en la hoja de datos del AD8004, que son los que proporcionan mayor estabilidad.
\end{itemize}



\section{Selección de componentes}
\label{sec:seleccion_componentes_deteccion}


% PLACEHOLDER: completar la tabla de componentes definitiva (valores de
% R4, C1, R7 y resistencias/capacitores de la red de realimentacion) una
% vez cerrado el diseno final de la placa comercial.

\section{Prototipado y validación experimental}
\label{sec:prototipado_validacion}


El circuito se modeló utilizando el modelo del amplificador AD8004 y se simuló en NGSPICE para distintos valores de tensión de entrada $V_{in}$. A partir de 25 barridos con diferentes valores de $V_{in}$ se verificó la linealidad del circuito a partir de $V_i = 0.8$~V, umbral fijado por la tensión de salida mínima que alcanzan los CFA en un esquema de alimentación de riel único de +12~V. La salida del detector se leyó 5~\textmu s después del pulso, instante en el que un ADC que muestrea a aproximadamente 500~ksps captura correctamente el valor retenido \cite{galvagno2026ilrc}.

Se ensambló un primer prototipo del circuito mediante la técnica de papel transfer y ataque químico con cloruro férrico, empleando
componentes de montaje superficial (SMD) dada la sensibilidad del circuito a las capacidades parásitas \cite{informeMarzo2026}. Al excitar
el circuito con un pulso de prueba de 100~ns de ancho y 20~Hz de frecuencia se observaron oscilaciones en los amplificadores
operacionales, atribuibles a las capacidades parásitas propias de la
condición de prototipo de la placa, que degradaban la linealidad
requerida del detector \cite{informeMarzo2026}.

Tras la fabricación industrial de una nueva placa (véase la
Sección~\ref{sec:mitigacion_parasitos}) se repitió el ensayo con un
pulso de 100~ns y se comprobó la eliminación de las oscilaciones
observadas previamente \cite{informeAbril2026}. La linealidad del
detector se verificó mediante 20 mediciones con saltos discretos de
amplitud de entrada de 200~mV, entre 1.2~V y 5~V, leyendo el valor de
salida tras un tiempo de asentamiento de 700~ns; los resultados fueron
consistentes con los obtenidos en la simulación \cite{informeAbril2026}.

% PLACEHOLDER: agregar comparacion cuantitativa entre los valores
% simulados y medidos (error relativo, linealidad, ancho de banda
% efectivo) cuando se disponga del analisis de datos completo.

\section{Diseño de PCB y mitigación de efectos parásitos}
\label{sec:mitigacion_parasitos}

Para reducir el efecto de las capacidades parásitas identificadas en la
etapa de prototipado, se rediseñó la placa aplicando las siguientes
buenas prácticas de ruteo \cite{informeMarzo2026}:

\begin{itemize}
  \item Mantenimiento de planos de masa con recortes (\emph{cutouts})
        debajo de los pines de entrada y de la red de realimentación.
  \item Ubicación de los componentes de la red de realimentación lo más
        cerca posible del amplificador, con el resistor $R_f$ conectado
        directamente entre el pin de salida y el pin $-IN$ sin vías
        intermedias.
  \item Colocación de capacitores de desacople en los pines de
        alimentación, con vía directa al plano de masa.
  \item Separación de las pistas de entrada respecto de líneas de alta
        frecuencia o señales ruidosas adyacentes.
  \item Reducción del encapsulado de las resistencias SMD de 1206 a
        0603, disminuyendo además su tolerancia de 5\,\% a 1\,\%.
\end{itemize}

La placa resultante se diseñó en KiCad~9.0
\cite{informeMarzo2026,kicad_hsrl} y se encargó inicialmente a la empresa
PCBWay y, posteriormente, a la empresa JLCPCB
\cite{informeMarzo2026,informeAbril2026}.